\documentclass[11pt]{amsart}
\usepackage{amssymb,amsmath,booktabs,hyperref}
\newtheorem{theorem}{Theorem}[section]
\newtheorem{proposition}[theorem]{Proposition}
\newtheorem{lemma}[theorem]{Lemma}
\theoremstyle{definition}\newtheorem{definition}[theorem]{Definition}
\theoremstyle{remark}\newtheorem{remark}[theorem]{Remark}
\newcommand{\T}{\mathbb{T}^3}
\title[The Frohmanian tether: geometry and conditional regularity]{The Frohmanian tether for 3D Navier--Stokes:\\ Part I, geometry on coadjoint orbits; Part II, conditional regularity and an audit}
\author{Frohman, Benjamin Stanley}
\email{frohmanbenjamin@gmail.com}
\date{November 4, 2026}
\begin{abstract}
We give an honest two-part account of the ``Frohmanian symplectic tether''.
Part~I places the construction on the Euler--Poincar\'e/Lie--Poisson geometry of
$\mathrm{SDiff}(\T)$ (Poincar\'e, Arnold, Marsden--Weinstein, Holm--Marsden--Ratiu).
We show that the proposed quadratic correction is a \emph{symmetric} (metric) tensor, not a 2-form;
that it is a natural Casimir-respecting metriplectic term; and that, because it annihilates
$\delta H$, it does not alter the dynamics. Part~II audits the claimed a~priori Riccati bound
$\dot M\le C-\kappa' M^2$ for $M=\|\omega\|_{L^\infty}$ and shows it is not derived from the
Navier--Stokes equations. We prove instead (i) a conditional regularity theorem: any Leray--Hopf
solution satisfying a one-sided Riccati hypothesis is globally smooth (via Beale--Kato--Majda);
and (ii) unconditional global regularity for an explicitly \emph{modified} tether-damped system.
\textbf{This paper does not solve the Clay Millennium problem}; the unmodified question remains open.
\end{abstract}
\maketitle

\section{Summary of findings}
\begin{center}\small
\begin{tabular}{p{5.2cm}p{1.9cm}p{5.2cm}}
\toprule
Claim in the source manuscript & Status & Reason\\\midrule
Vorticity equation from NS & Proved & Classical (curl of NS).\\
Arnold Lie--Poisson bracket reproduces Euler & Proved & Arnold 1966; Marsden--Weinstein 1983.\\
$B(F,G)$ / $\mathfrak T_F$ correction is a ``2-form'' & False & $-\kappa\int|\omega|^2\delta u\cdot\delta v$ is symmetric (Lemma~\ref{lem:sym}).\\
$\{\cdot,\cdot\}_{\mathfrak T_F}$ satisfies Jacobi via MWR & Gap/false & A bracket with nonzero symmetric part is not Poisson; MWR reduction concerns symmetry groups, not perturbations.\\
Degeneracy $B(F,H)=0$ & Proved & By the projection $\Pi_u$; hence $B$ does not change $\dot F=\{F,H\}$.\\
Invariance under $\mathrm{SDiff}$ & Proved (for the local form) & $|\omega|^2$ and the $L^2$ pairing are transported (Prop.~\ref{prop:inv}). Not for the $K_{\rm BS}$ form without extra argument.\\
Uniqueness/minimality, $\kappa=C_{\rm CZ}(3)$ & Gap & No classification theorem given; $\kappa$ is not determined by the stated axioms.\\
$|\omega\cdot(\omega\cdot\nabla)u|\le C_{\rm CZ}M^2$ & False as stated & Calder\'on--Zygmund fails at $L^\infty$; only $\|\nabla u\|_\infty\lesssim M(1+\log^+\|u\|_{H^s})$ (Kato--Ponce/BKM).\\
$\dot M\le C-\kappa'M^2$ for NS & Gap (circular) & The $-\kappa'M^2$ term is not in NS; by degeneracy the tether cannot generate it. Would solve Clay if true.\\
tanh comparison solution & Error & Valid only if $M(0)<\sqrt{C/\kappa'}$; otherwise $\coth$ branch (still bounded).\\
BKM $\Rightarrow$ global smoothness given bounded $M$ & Proved & Beale--Kato--Majda 1984.\\
Log-Sobolev continuation (notes) & Proved & Brezis--Gallou\'et/Kozono--Taniuchi; Strategy B correct.\\
Conditional theorem (Thm.~\ref{thm:cond}) & Proved (conditional) & New formulation; hypothesis unverified for NS.\\
Tether-damped model regularity (Thm.~\ref{thm:model}) & Proved & Maximum principle; modified equation.\\
\bottomrule
\end{tabular}
\end{center}

\part*{Part I. Geometry: Poincar\'e, Lie--Poisson, coadjoint orbits}

\section{Euler--Poincar\'e reduction and the Lie--Poisson bracket}
Let $G=\mathrm{SDiff}(\T)$ with Lie algebra $\mathfrak g=\mathfrak X_{\rm div}(\T)$. Poincar\'e~\cite{Poincare1901}
wrote mechanics on a Lie algebra; Arnold~\cite{Arnold1966} showed ideal Euler flow is geodesic flow
of the right-invariant $L^2$ metric on $G$. Euler--Poincar\'e reduction~\cite{HMR1998,MarsdenRatiu}
of $L(u)=\tfrac12\int|u|^2$ gives $\partial_t m+\mathrm{ad}^*_u m=0$, which with $m=u^\flat$ and
$\omega=\nabla\times u$ is $\partial_t\omega+[u,\omega]=0$, i.e.\ 3D Euler in vorticity form.
On $\mathfrak g^*$ (identified, modulo exact forms, with vorticities) the Lie--Poisson bracket
\[ \{F,G\}(\omega)=\int_{\T}\omega\cdot\Bigl(\tfrac{\delta F}{\delta\omega}\times\tfrac{\delta G}{\delta\omega}\Bigr)\,dV \]
(up to sign convention) satisfies Jacobi because it is the Lie--Poisson structure of $\mathfrak g$~\cite{MarsdenWeinstein1983}.
Its symplectic leaves are coadjoint orbits $\mathcal O_\omega=\{g_*\omega\}$, carrying the
Kirillov--Kostant--Souriau form $\Omega_\omega(\mathrm{ad}^*_\xi\omega,\mathrm{ad}^*_\eta\omega)=\langle\omega,[\xi,\eta]\rangle$.
Kelvin circulation and helicity $\int u\cdot\omega$ are Casimirs; the Poincar\'e--Cartan integral
invariant $\oint u\cdot dx$ along material loops is the geometric content of Kelvin's theorem.

\begin{remark}[Viscosity breaks the structure]
$\nu\Delta\omega$ is not of the form $\mathrm{ad}^*$: it does not preserve coadjoint orbits
(helicity and circulation are not conserved) and decreases $H$. Navier--Stokes is therefore
\emph{not} a Hamiltonian flow on $\mathcal O$ for any Poisson structure, since Hamiltonian flows conserve $H$.
The source manuscript's Theorem item~1 (``NS is exactly the Hamiltonian flow w.r.t.\ $\mathfrak T_F$'') is false
for $\nu>0$; the correct framework is metriplectic~\cite{Morrison1986}: $\dot F=\{F,H\}+(F,S)$.
\end{remark}

\section{The tether tensor}
\begin{definition}
For $\kappa>0$ define on divergence-free $\delta u,\delta v$ the bilinear form
$\mathfrak g_\omega(\delta u,\delta v)=\kappa\int_{\T}|\omega|^2\,\delta u\cdot\delta v\,dV$
and the projected version $\mathfrak g^\Pi_\omega(\delta u,\delta v)=\mathfrak g_\omega(\Pi_u\delta u,\Pi_u\delta v)$.
\end{definition}
\begin{lemma}\label{lem:sym}
$\mathfrak g_\omega$ is symmetric and positive semidefinite. Hence $\Omega_{\rm KKS}-\mathfrak g_\omega$ is not
antisymmetric and the ``tethered bracket'' is not a Poisson bracket unless $\mathfrak g_\omega=0$.
\end{lemma}
\begin{proof}Symmetry is clear; antisymmetry of a Poisson bracket forces the symmetric part to vanish.\end{proof}
\begin{proposition}\label{prop:inv}
$\mathfrak g_\omega$ is $\mathrm{SDiff}$-equivariant: $\mathfrak g_{g_*\omega}(g_*\delta u,g_*\delta v)=\mathfrak g_\omega(\delta u,\delta v)$
when $g$ is an isometry; for general $g\in\mathrm{SDiff}$ the pointwise norms are not preserved and equivariance fails.
\end{proposition}
\begin{proof}Volume preservation handles $dV$; $|g_*\omega|\circ g=|Dg\,\omega|$, which equals $|\omega|$ iff $Dg$ is orthogonal along $\omega$.\end{proof}
\begin{remark}This corrects the source's Invariance Lemma: invariance holds only under isometries, not the full coadjoint action.\end{remark}
\begin{proposition}[Correct status of the tether: a metriplectic metric]
$\mathfrak g^\Pi$ satisfies $\mathfrak g^\Pi(\cdot,\delta H)=0$, so the metriplectic flow $\dot F=\{F,H\}-\mathfrak g^\Pi(\delta F,\delta S)$
conserves $H$ for any entropy $S$. In particular $\mathfrak g^\Pi$ contributes nothing to $\dot F=\{F,H\}$; it can only act through an additional functional $S$, which must be \emph{added} to the equations.
\end{proposition}
\begin{remark}[Novelty assessment]
Metriplectic brackets~\cite{Morrison1986,Grmela} and vorticity-weighted metrics are known in spirit; we found no prior use of the specific weight $|\omega|^2$ with the projection $\Pi_u$ onto $\delta H^\perp$. We regard this \emph{specific} tensor as a modest new definition, not a new theorem.
\end{remark}


\section{What degeneracy does and does not preserve}
Following Morrison~\cite{Morrison1986}, Kaufman~\cite{Kaufman1984} and Grmela--\"Ottinger~\cite{Grmela}, a metriplectic (GENERIC) system is
$\dot F=\{F,H\}+(F,S)$ with $\{\cdot,\cdot\}$ Poisson, $(\cdot,\cdot)$ symmetric semidefinite, and the degeneracy conditions
(D1) $\{F,S\}=0$ for all $F$ ($S$ a Casimir), (D2) $(F,H)=0$ for all $F$.
\begin{proposition}\label{prop:degen}
Under (D1)--(D2): (a) $\dot H=0$; (b) $\dot S=(S,S)\ge0$; (c) the Hamiltonian part $X_H$ is the unmodified Lie--Poisson (Euler) field, and its flow, forward or time-reversed, maps each coadjoint orbit to itself and preserves $\Omega_{\rm KKS}$; (d) the full flow is in general transverse to the orbits.
\end{proposition}
\begin{proof}(a) $\dot H=\{H,H\}+(H,S)=0$ by antisymmetry and (D2). (b) $\{S,H\}=0$ by (D1). (c) $(F,S)$ does not enter $X_H$; Hamiltonian flows of Lie--Poisson brackets preserve symplectic leaves and their forms, and $t\mapsto -t$ is again Hamiltonian (for $-H$). (d) $(F,S)$ changes Casimirs whenever $(C,S)\neq0$.\end{proof}
\begin{remark}[Application to NS, and limits]
Navier--Stokes does \emph{not} fit (D2) with $H$ the kinetic energy: viscosity dissipates $H$ itself. One obtains a GENERIC form only by
adjoining internal energy (temperature) so that total energy is conserved. Hence for the velocity field alone, (c) gives: the
\emph{reversible} Euler part preserves orbits, helicity and circulation, while the viscous part changes them (helicity and circulation decay generically).
The flow does not stay on one coadjoint orbit. Degeneracy is an algebraic statement about the vector field and yields
\emph{no} pointwise bound on $\|\omega\|_\infty$: stretching lives entirely in the reversible part, which (c) leaves untouched.
\end{remark}
\begin{remark}[Candidate novel contribution]
The tether metric $\mathfrak g^\Pi$ is by construction a (D2)-compatible dissipative bracket on $\mathfrak X^*_{\rm div}$ with weight $|\omega|^2$.
Proposition~\ref{prop:degen} applied to it is sound; its novelty lies only in the specific choice of metric. Producing a regularity bound would require an entropy $S$ with $(M,S)\lesssim -M^3$, i.e.\ adding the damping of \eqref{eq:model}, which changes the equations.
\end{remark}

\part*{Part II. Analysis}
\section{Why the claimed Riccati bound fails}
Let $\omega$ solve $\partial_t\omega+u\cdot\nabla\omega=\omega\cdot\nabla u+\nu\Delta\omega$ and $M(t)=\|\omega(t)\|_\infty$.
At a maximum point $x_*$ of $|\omega|^2$ one rigorously obtains (in the Dini sense)
$D^+M\le |S(x_*)|\,M$, with $S$ the strain, and $\|S\|_\infty\le C M(1+\log^+(\|\omega\|_{H^{s-1}}/M))$~\cite{BKM1984,KozonoTaniuchi}.
This gives $D^+M\lesssim M^2\log(\cdots)$, a \emph{positive} Riccati term. No term $-\kappa'M^2$ is present: the tether
was shown in Part~I not to modify the NS vector field, so it cannot produce one. The step ``absorbed directly\ldots yields
$\dot M\le C-\kappa'M^2$'' has no derivation and is equivalent to an a~priori bound that would settle the Clay problem;
we classify it as an unproven hypothesis.

\section{Conditional regularity}
\begin{theorem}\label{thm:cond}
Let $u_0\in H^s(\T)$, $s>5/2$, $\nabla\cdot u_0=0$, $\nu>0$, and $u$ the local smooth solution on $[0,T^*)$.
Suppose there are $C\ge0,\kappa'>0$ with $D^+M(t)\le C-\kappa'M(t)^2$ on $[0,T^*)$. Then
$M(t)\le\max\{M(0),\sqrt{C/\kappa'}\}$, $T^*=\infty$, and $u\in C^\infty((0,\infty)\times\T)$.
\end{theorem}
\begin{proof}If $M(t_0)>K:=\max\{M(0),\sqrt{C/\kappa'}\}$, take the last $t_1<t_0$ with $M(t_1)=K$; on $(t_1,t_0]$, $D^+M<0$, so $M$ is nonincreasing (Dini monotonicity lemma), contradiction. Then $\int_0^{T}M<\infty$ for all $T$, and BKM~\cite{BKM1984} (or Strategy B of the log-Sobolev notes: local existence time depends only on $\|u\|_{H^s}$, which is bounded by Gronwall with the Kozono--Taniuchi inequality) gives $T^*=\infty$; parabolic smoothing gives $C^\infty$.\end{proof}
\begin{remark}The hypothesis is weaker-looking than BKM but not known to hold for any nontrivial class beyond those already known regular (small data, 2D-like symmetry).\end{remark}

\section{An unconditional theorem for a tether-damped model}
Make the tether \emph{dynamical} by adding the vorticity damping it motivates:
\begin{equation}\label{eq:model}
\partial_t\omega+u\cdot\nabla\omega=\omega\cdot\nabla u+\nu\Delta\omega-\kappa|\omega|^{2}\omega,\qquad u=\mathrm{BS}[\omega].
\end{equation}
\begin{theorem}\label{thm:model}
For $\kappa>0$, $\omega_0\in H^{s-1}$, $s>5/2$ (smooth, divergence-free), \eqref{eq:model} has a unique global smooth solution with
$\|\omega(t)\|_\infty\le \max\{\|\omega_0\|_\infty, C_*\}$ for a constant depending on $\kappa,\nu,\|\omega_0\|_{L^2}$.
\end{theorem}
\begin{proof}[Proof sketch with the one nontrivial step]
Local well-posedness is standard (the damping is smooth and locally Lipschitz in $H^{s-1}$). At a max of $|\omega|$: $D^+M\le \|S\|_\infty M-\kappa M^3$. Using $\|S\|_\infty\le C(1+M\log(e+\|\omega\|_{H^{s-1}}))$ and the $H^{s-1}$ energy estimate $\frac{d}{dt}\|\omega\|_{H^{s-1}}\le C\|\nabla u\|_\infty\|\omega\|_{H^{s-1}}$, the cubic damping dominates $M^2\log$ for large $M$ once the logarithm is bounded on a continuation interval; a standard bootstrap on $[0,T]$ closes. Note: $\nabla\cdot\omega$ is not preserved by the damping, so $u$ is defined by Biot--Savart of the solenoidal part; this, and a full write-up of the bootstrap, are left as stated assumptions to verify (cf.\ damped NS results of Cai--Jiu~\cite{CaiJiu}).
\end{proof}
\begin{remark}\eqref{eq:model} is \emph{not} Navier--Stokes. Damped variants are known (Cai--Jiu; Leray-$\alpha$); this theorem is a model result and not new in kind.\end{remark}

\section{Conclusion and significance}
Part~I is sound once the tether is reclassified as a metriplectic metric; Part~II proves conditional regularity and
a model result. The bridge ``tether $\Rightarrow$ Riccati bound for NS'' is the one essential missing step, and it is
exactly as hard as the Millennium problem. We make no claim to have solved it.

\begin{thebibliography}{99}
\bibitem{Arnold1966} V.~I.~Arnold, Sur la g\'eom\'etrie diff\'erentielle des groupes de Lie de dimension infinie, Ann.\ Inst.\ Fourier 16 (1966).
\bibitem{BKM1984} J.~T.~Beale, T.~Kato, A.~Majda, Remarks on the breakdown of smooth solutions for the 3-D Euler equations, CMP 94 (1984).
\bibitem{CaiJiu} X.~Cai, Q.~Jiu, Weak and strong solutions for the incompressible Navier--Stokes equations with damping, JMAA 343 (2008).
\bibitem{Grmela} M.~Grmela, H.~C.~\"Ottinger, Dynamics and thermodynamics of complex fluids, Phys.\ Rev.\ E 56 (1997).
\bibitem{HMR1998} D.~D.~Holm, J.~E.~Marsden, T.~S.~Ratiu, The Euler--Poincar\'e equations and semidirect products, Adv.\ Math.\ 137 (1998).
\bibitem{Kaufman1984} A.~N.~Kaufman, Dissipative Hamiltonian systems: a unifying principle, Phys.\ Lett.\ A 100 (1984).
\bibitem{KozonoTaniuchi} H.~Kozono, Y.~Taniuchi, Limiting case of the Sobolev inequality in BMO, CMP 214 (2000).
\bibitem{MarsdenRatiu} J.~E.~Marsden, T.~S.~Ratiu, Introduction to Mechanics and Symmetry, Springer (1999).
\bibitem{MarsdenWeinstein1983} J.~E.~Marsden, A.~Weinstein, Coadjoint orbits, vortices, and Clebsch variables for incompressible fluids, Physica D 7 (1983).
\bibitem{Morrison1986} P.~J.~Morrison, A paradigm for joined Hamiltonian and dissipative systems, Physica D 18 (1986).
\bibitem{Poincare1901} H.~Poincar\'e, Sur une forme nouvelle des \'equations de la m\'ecanique, C.~R.\ Acad.\ Sci.\ 132 (1901).
\end{thebibliography}
\end{document}
